\section{Fazit}
\label{sec:conclusion}

Dieses Projekt hat untersucht, ob Hessian-basierte Krümmungsmetriken vorhersagen können, welche Schichten eines CNN unter der Power-of-Two-Quantisierung am stärksten geschädigt werden, und liefert auf alle drei Forschungsfragen eine klare, wenn auch nuancierte Antwort.
\par\medskip

Zu \textbf{F1} (schichtweise Krümmungsstruktur unter PTQ vs.\ QAT): Die Random-Init-Kontrolle (\S\ref{sec:phase2}) zeigt, dass die für die Analyse zentrale \texttt{conv1}-Krümmungsspitze in allen drei Architekturen durch Training entsteht und nicht bereits bei Initialisierung vorhanden ist; die Fusions-Analyse (\S\ref{sec:phase35}) zeigt weiter, dass diese Spitze überwiegend ein Effekt der Batch-Normalisierungs-Fusion ist und nicht der PoT-Quantisierung selbst, mit Ausnahme des \texttt{cnn}, wo PTQ zusätzlich moderat verstärkt. Der direkte PTQ-vs-QAT-Vergleich (\S\ref{sec:phase9}) zeigt ein zunächst kontraintuitives Bild: Die relative Krümmungsstruktur selbst bleibt zwischen PTQ- und QAT-Checkpoints nahezu unverändert ($\rho=0.95$--$1.00$ über alle sechs Kombinationen), während die Beziehung dieser Krümmung zum tatsächlichen Schaden unter QAT fast vollständig zusammenbricht (0 von 18 signifikant, gegenüber 5 von 18 unter PTQ). Der Grund liegt nicht in der Krümmung, sondern im Ziel. Die isolierte Schichtquantisierung verbessert unter QAT häufig den Loss gegenüber der FP32-Baseline, statt ihn zu verschlechtern (\S\ref{sec:phase1-results}), sodass dieselbe Krümmungslandschaft ein grundlegend anderes Zielsignal vorherzusagen versucht (\S\ref{sec:disc-ptqqat}).
\par\medskip

Zu \textbf{F2} (der Kernfragestellung): Kein einzelner Score ist universell zuverlässig, um die am stärksten geschädigte Schicht zu identifizieren. Die rohe Krümmung ($S_{\mathrm{raw}}$) identifiziert die wahre Top-Damage-Schicht in 3 von 6 Kombinationen exakt (Rang~1) und wird in ihrer top-$k^*$-Überlappung von keinem der beiden anderen Scores übertroffen. Das HAWQ-V2-Produkt scheitert spezifisch daran, dass die PoT-Störungsgröße mit der wahren Sensitivität \emph{anti-korreliert}, statt neutral oder positiv zu sein (\S\ref{sec:disc-hawq}), was ein PoT-spezifischer Mechanismus und keine allgemeine Schwäche Krümmungs-basierter Sensitivität ist. Welche Schicht konkret am stärksten geschädigt wird, hängt zudem stark von Architektur und Datensatz ab (\texttt{conv1} auf ImageNet100 in allen drei Architekturen, aber uneinheitlich auf CIFAR10, Tabelle~\ref{tab:rank}).
\par\medskip

Zu \textbf{F3} (Gewichts- vs.\ Aktivierungsregime und Anwendbarkeit): Unter PTQ auf CIFAR10 dominiert der Gewichtsschaden in allen drei Architekturen, unter QAT dominiert der Aktivierungsschaden (\S\ref{sec:phase1-results}). Die Hessian-basierte Sensitivitätsanalyse, die ausschließlich Gewichtsschaden erfasst, ist damit nicht universell anwendbar, sondern regimeabhängig, und ihre Anwendbarkeit sollte vor jedem Einsatz explizit geprüft werden.
\par\medskip

Über diese drei Fragen hinaus liefert das Projekt zwei methodische Beiträge mit Reichweite über den eigenen Anwendungsfall hinaus: das Rekonziliations-Ledger als dokumentierter Beleg dafür, wie stark Trace-Messungen von unausgesprochenen Konfigurationsentscheidungen abhängen (\S\ref{sec:phase3}), und die Deployment-Evaluation als Beleg dafür, dass ein theoretischer Hardware-Vorteil ohne ausgereifte Software-Unterstützung nicht automatisch real wird (\S\ref{sec:phase10}). Ebenso gehört das dokumentierte Scheitern der KFAC-Zerlegung (\S\ref{sec:phase6}) zu den ehrlich zu berichtenden Ergebnissen dieses Projekts.
\par\medskip

Die Kernergebnisse zu F2, also der direkte Test der HAWQ-V2-Formulierung gegen gemessenen Schaden, das Rekonziliations-Ledger und die Metrik-Konkordanz-Analyse, wurden in kompakter Form als Workshop-Paper bei \textbf{AXIOM @ NeurIPS 2026} eingereicht. Dieser Bericht dokumentiert das Projekt in seiner vollen Breite, einschließlich der Phasen, Kontrollen und Negativresultate, die im begrenzten Format eines Workshop-Papers keinen Platz fanden, aber für ein vollständiges Bild der methodischen Arbeit dieses Projekts wesentlich sind.